\documentclass[11pt,a4paper]{article}

\usepackage{fontspec}
\setmainfont{DejaVu Sans}
\setsansfont{DejaVu Sans}
\setmonofont{DejaVu Sans Mono}[Scale=0.88]
\defaultfontfeatures{Ligatures=TeX}
\usepackage{microtype}
\usepackage{unicode-math}
\setmathfont{Latin Modern Math}[Scale=1.0]
\usepackage[ngerman]{babel}
\usepackage{geometry}
\geometry{margin=2.5cm}
\usepackage{amsmath,amssymb}
\usepackage{parskip}
\usepackage{booktabs}
\usepackage{array}
\usepackage{xcolor}
\usepackage{hyperref}
\hypersetup{
  colorlinks=true,
  linkcolor=blue!50!black,
  urlcolor=blue!50!black,
  pdftitle={Die Casimir-Skala in FFGFT - Kurzfassung},
  pdfauthor={Johann Pascher}
}

\title{Die Casimir-Skala in FFGFT\\[2pt]
  \large Kurzfassung für die IPI-Korrespondenz}
\author{Johann Pascher\\[2pt]
\small ORCID: \href{https://orcid.org/0009-0000-6518-4064}{0009-0000-6518-4064}\\
\small GitHub: \url{https://github.com/jpascher/T0-Time-Mass-Duality}\\
\small Zenodo (FFGFT v1.0.13): \url{https://zenodo.org/records/20041529}\\
\small DOI: \href{https://doi.org/10.5281/zenodo.20041529}{10.5281/zenodo.20041529}}
\date{Mai 2026}

\begin{document}
\maketitle

\begin{abstract}
\noindent Die Fundamentale Fraktale Geometrische Feldtheorie (FFGFT) macht eine quantitative, parameterfreie Vorhersage für die charakteristische Casimir-Vakuumskala: $L_\xi \approx 100{,}24\,\mu$m. Die Vorhersage wird aus der Energiedichte der kosmischen Hintergrundstrahlung (CMB) und einer einzigen dimensionslosen Konstante $\xi = 4/30{,}000$ abgeleitet. Diese Notiz fasst die Herleitung zusammen, die Konsistenz mit der Standard-Casimir-Formel und die Einbettung von $L_\xi$ in eine fünfstufige $\xi$-Potenzleiter, die die Planck-Länge mit der Hubble-Länge verbindet. Quelldokumente im FFGFT-Korpus: Doc~091 (Casimir-CMB), Doc~165 (Hubble), Doc~166 (Synthese).
\end{abstract}

\section{Ausgangspunkt}

FFGFT arbeitet mit einem einzigen dimensionslosen Parameter,
\[
\xi = \frac{4}{30{,}000} = 1{,}333 \times 10^{-4},
\]
und der Skalen-Anker-Beziehung $\tilde T \cdot m = 1$. Aus diesen und Standard-Naturkonstanten leitet FFGFT eine Hierarchie physikalischer Längenskalen ohne weitere Eingabeparameter ab. Zwei dieser Skalen sind für den Casimir-Effekt direkt relevant: eine Sub-Planck-Sättigungslänge $L_0 = \xi \cdot \ell_P$ und eine vakuum-emergente Länge $L_\xi$, verknüpft mit der CMB.

\section{Die CMB-Casimir-Beziehung}

Die zentrale Beziehung in Doc~091 lautet:
\begin{equation}
\rho_{\mathrm{CMB}} = \frac{\xi \, \hbar c}{L_\xi^4}\,.
\label{eq:cmb_casimir}
\end{equation}
Mit der gemessenen CMB-Energiedichte $\rho_{\mathrm{CMB}} = 4{,}17 \times 10^{-14}$ J/m$^3$ und $\xi = 4/30{,}000$ legt die Beziehung fest:
\begin{equation}
L_\xi = \left(\frac{\xi\,\hbar c}{\rho_{\mathrm{CMB}}}\right)^{1/4} \approx 100{,}24\,\mu\text{m}.
\label{eq:Lxi_value}
\end{equation}
Äquivalent über die Stefan-Boltzmann-Form $\rho_{\mathrm{CMB}} = (\pi^2/15)(k_B T_{\mathrm{CMB}})^4 / (\hbar c)^3$:
\begin{equation}
L_\xi = \left(\frac{15\,\xi}{\pi^2}\right)^{1/4} \frac{\hbar c}{k_B T_{\mathrm{CMB}}} = 100{,}24\,\mu\text{m}\,,
\label{eq:Lxi_from_T}
\end{equation}
in Übereinstimmung mit der direkten Auswertung auf $0{,}11\%$.

Die zugehörige charakteristische Frequenz ist
\[
f_\xi = \frac{c}{L_\xi} \approx 2{,}99\,\text{THz},
\]
was $L_\xi$ in dasselbe logarithmische Regime stellt wie die in der IPI-Korrespondenz diskutierten Casimir-Zielfrequenzen (Moseleys $\nu_M \approx 8{,}23$~THz / $\lambda_M \approx 36{,}4\,\mu$m), mit einer Trennung um Faktor $\sim 2{,}7$.

\section{Konsistenz mit der Standard-Casimir-Formel}

Die modifizierte Casimir-Energiedichte in FFGFT lautet
\begin{equation}
|\rho_{\mathrm{Casimir}}(d)| = \frac{\pi^2}{240\,\xi}\,\rho_{\mathrm{CMB}}\left(\frac{L_\xi}{d}\right)^4\,.
\label{eq:modified_casimir}
\end{equation}
Einsetzen von Gl.~\eqref{eq:cmb_casimir} in Gl.~\eqref{eq:modified_casimir} liefert
\begin{equation}
|\rho_{\mathrm{Casimir}}(d)| = \frac{\pi^2}{240\,\xi}\cdot\frac{\xi\,\hbar c}{L_\xi^4}\cdot\frac{L_\xi^4}{d^4} = \frac{\pi^2 \hbar c}{240\,d^4}\,,
\end{equation}
exakt die Standard-Casimir-Formel. Der Parameter $\xi$ kürzt sich auf operationaler Ebene heraus. Die FFGFT-Formulierung ist daher konsistent mit allen aktuellen Casimir-Messungen, die die Standard-$1/d^4$-Skalierung bestätigen, und reformuliert dabei die absolute Energieskala über $L_\xi$ und $\rho_{\mathrm{CMB}}$.

\section{Die Längenskalen-Hierarchie}

FFGFT verortet $L_\xi$ auf einer fünfstufigen Leiter physikalischer Längenskalen, alle ausdrückbar als $\xi^n \cdot \ell_P$:

\begin{table}[h]
\centering
\small
\begin{tabular}{lccc}
\toprule
\textbf{Skala} & \textbf{Bedeutung} & \textbf{Länge [m]} & $\xi^n \cdot \ell_P$ \\
\midrule
$L_0 = \xi \cdot \ell_P$ & Sub-Planck-Sättigung & $2{,}16 \times 10^{-39}$ & $n = +1{,}0$ \\
$\ell_P$ & Planck-Länge & $1{,}62 \times 10^{-35}$ & $n = 0$ \\
$\bar\lambda_e$ & Compton-Wellenlänge des Elektrons & $3{,}86 \times 10^{-13}$ & $n = -5{,}77$ \\
$L_\xi$ & Casimir-Vakuumskala & $1{,}00 \times 10^{-4}$ & $n = -7{,}95$ \\
$c/H_0$ & Hubble-Länge & $1{,}38 \times 10^{26}$ & $n = -15{,}72$ \\
\bottomrule
\end{tabular}
\caption{Die $\xi$-Potenzleiter spannt sich von Sub-Planck- bis zu kosmologischen Skalen. $L_\xi$ sitzt auf einer regelmäßigen Leiter, nicht als isolierter freier Parameter.}
\end{table}

Eine Beobachtung zum geometrischen Mittel (Doc~166):
\begin{equation}
\sqrt{\ell_P \cdot \tfrac{c}{H_0}} = 47{,}3\,\mu\text{m}, \qquad L_\xi = 100{,}2\,\mu\text{m}, \qquad \frac{L_\xi}{\sqrt{\ell_P \cdot c/H_0}} = 2{,}12\,.
\end{equation}

\section{Die Casimir-Hubble-Brückenrelation}

FFGFT sagt auch eine algebraische Beziehung zwischen der Casimir-Skala und der Hubble-Konstante voraus. Kombiniert man Gl.~\eqref{eq:Lxi_from_T} mit der FFGFT-Vorhersage für $H_0$ in Doc~165, $\hbar H_0 / (m_e c^2) = (\pi/2)\,\xi^{10}$, ergibt sich:
\begin{equation}
\boxed{\;L_\xi \cdot \frac{H_0}{c} = \frac{\pi}{2}\left(\frac{15}{\pi^2}\right)^{1/4} \frac{m_e c^2}{k_B T_{\mathrm{CMB}}}\,\xi^{41/4}\;}
\label{eq:bridge}
\end{equation}
Der Exponent zerlegt sich als
\[
\frac{41}{4} = \underbrace{10}_{H_0\text{-Exponent (Doc 165)}} + \underbrace{\frac{1}{4}}_{\text{Casimir-Exponent (Doc 091)}}\,,
\]
derselbe Exponent, der in Doc~026 als $E_H = E_0 \cdot \xi^{41/4}$ erscheint. Numerisch stimmen beide Seiten von Gl.~\eqref{eq:bridge} bis $10^{-14}\%$ überein. Die Beziehung enthält keine freien Parameter; alle Größen sind entweder Naturkonstanten oder das einzelne $\xi$.

In numerischer Form:
\begin{equation}
L_\xi \cdot \frac{H_0}{c} = 7{,}241 \times 10^{-31}\,, \qquad \frac{m_e c^2}{k_B T_{\mathrm{CMB}}} = 2{,}000 \times 10^9\,.
\end{equation}

\section{Vergleich mit experimentellen Casimir-Daten}

Berichtete experimentelle Casimir-Messungen effektiver $L_\xi$-Werte (aus Doc~091):

\begin{table}[h]
\centering
\small
\begin{tabular}{lcc}
\toprule
\textbf{Konfiguration} & \textbf{Effektives $L_\xi$} & \textbf{Referenz} \\
\midrule
parallele Platten & $228$ nm & Dhital 2024 \\
Kugel-Platte & $1{,}75\,\mu$m & Xu 2022 \\
weiterer Wert & $18\,\mu$m & verschiedene \\
\bottomrule
\end{tabular}
\caption{Berichtete experimentelle $L_\xi$-Werte streuen über Konfigurationen (228~nm bis 18~$\mu$m), spiegelnd geometrische Unterschiede ($F \propto 1/L^4$ für parallele Platten, $F \propto 1/L^3$ für Kugel-Platte) und experimentelle Bedingungen. Die FFGFT-Vorhersage $L_\xi \approx 100\,\mu$m ist der konfigurationsunabhängige Vakuum-CMB-Anker; geometriespezifische Werte weichen erwartungsgemäß um bekannte Faktoren ab.}
\end{table}

\section{Zusammenfassung der Vorhersagen}

\begin{enumerate}\itemsep2pt
\item $L_\xi = 100{,}24\,\mu$m \emph{folgt} aus $\rho_{\mathrm{CMB}}$ und einem einzigen dimensionslosen $\xi = 4/30{,}000$. Keine freien Anpassungsparameter.
\item Einsetzen in die modifizierte Casimir-Dichte Gl.~\eqref{eq:modified_casimir} reproduziert das Standard-$\pi^2 \hbar c/(240\,d^4)$ exakt. Mathematische Konsistenz mit allen Standard-Casimir-Messungen auf operationaler Ebene bleibt erhalten.
\item $L_\xi$ sitzt auf einer fünfstufigen $\xi$-Potenzleiter von Planck- bis Hubble-Skala; sie ist näherungsweise das $2{,}12$-fache des geometrischen Mittels $\sqrt{\ell_P \cdot c/H_0}$.
\item Casimir-Skala und Hubble-Konstante erfüllen die algebraische Brückenrelation Gl.~\eqref{eq:bridge}, mit kombiniertem $\xi$-Exponenten $41/4 = 10 + 1/4$.
\item Die zugehörige Frequenz $c/L_\xi \approx 3$ THz liegt im selben Regime wie andernorts in der IPI-Korrespondenz diskutierte Casimir-Phänomenologie.
\end{enumerate}

\section*{Dokumentenverweise}

\noindent\begin{tabular}{ll}
\textbf{Doc 091} & \emph{Vereinheitlichte Casimir-CMB-Beschreibung.} Herleitung von $L_\xi$ und der \\
& modifizierten Casimir-Formel; Konsistenz mit der Standard-$1/d^4$-Skalierung. \\[2pt]
\textbf{Doc 165} & \emph{Das Hubble-Verhältnis.} Herleitung von $\hbar H_0 / (m_e c^2) = (\pi/2)\,\xi^{10}$, \\
& $H_0 \approx 66{,}8$ km/s/Mpc innerhalb 0,9\% des Planck-Wertes. \\[2pt]
\textbf{Doc 166} & \emph{Die Casimir-Skala als kosmisches Bindeglied.} Die $\xi$-Potenzleiter \\
& und die Brückenrelation $L_\xi \cdot H_0/c$; Erklärung des Exponenten $41/4$. \\[2pt]
\textbf{Doc 026} & \emph{Geometrische Kosmologie.} Der Exponent $41/4$ in $E_H = E_0 \cdot \xi^{41/4}$. \\
\end{tabular}

\medskip
Alle Dokumente verfügbar unter \url{https://github.com/jpascher/T0-Time-Mass-Duality}\\
und als Zenodo-Release v1.0.13: \url{https://zenodo.org/records/20041529}.

\end{document}
